\hspace{3mm}

\centerline{\bf \Huge Вступительный экзамен по алгебре}


\begin{flushright}
        \scriptsize{\textbf{Математика — это единственный совершенный метод водить самого себя за нос. Альберт Эйнштейн}}
\end{flushright}

\problem{\textbf{1}} Решите системы из двух линейных алгебраических уравнений (СЛАУ):

а.\textbf{(3)} $ \begin{cases}
   x + 2y = 25
   \\
   11x + 22y = 52
 \end{cases}$

б.\textbf{(3)} $\begin{cases}
    3x - 4y = 25
   \\
   5y + 11x = 1
 \end{cases}$

в.\textbf{(3)} $\begin{cases}
   -9x + 2y = 19
   \\
   63x - 14y = -133
 \end{cases}$


\problem{\textbf{2}}\textbf{(30)} Чему равно следующее выражение: 

$\Bigr(\dfrac{2^{14,54} \cdot 3^{8,77}}{12^{7,77}}\Bigl)\cdot\Bigr(\dfrac{(1,238+2,762) \cdot \dfrac{11}{110}}{(36,487-34,237): 2,8125}+\dfrac{(4,36-1,16) \cdot 0,3125}{\dfrac{9}{45} \cdot(47,8-45,55): 0,225}\Bigl)\cdot$ 

$\cdot \Bigr(39,86-\dfrac{8 \dfrac{4}{5}}{6,934+38 \dfrac{13}{40}: 2 \dfrac{1}{2} \cdot 0,2}-17\Bigl)\cdot\Bigr(7,5^2-2,5^2\Bigl)\cdot\Bigr(\dfrac{20\dfrac{8}{15}\cdot7,5-54,6:\dfrac{2}{5}}{3\dfrac{13}{21}\cdot8,4-34,4:14\dfrac{1}{3}}\Bigl)\cdot12,8--3188=?$



\problem{\textbf{3}} а.\textbf{(3)} Напишите все формулы сокращенного умножения (ФСУ), которые вы знаете. 

б.\textbf{(24)} Чему равно следующее выражение: 

$152,67^{3}-32,57^{3}+49992,5^{2}+458,01\cdot32,57^{2}-50007,5^{2}-152,67^{2}\cdot97,71=?$

\problem{\textbf{4}} а.\textbf{(12)} Постройте графики двух прямых: $y=4x-5$ и $y=-0,25x+3$, а также найдите координаты их общей точки.

б.\textbf{(24)} Прямая $A$ проходит через точки $B(1;-2)$ и $C(5;6)$. Прямая $D$ проходит через точку $E(-2;-3)$ и параллельна прямой $A$. Прямые $y=-x-3$ и $y=5x+9$ пересекаются в точке $F$. Прямая $G$ проходит через точку $F$ и перпендикулярна прямой $D$. Прямые $D$ и $G$ пересекаются в точке $H$. Найдите сумму абсциссы и ординаты точки $H$, умноженную на -10, и прибавьте к этому 2000.

\hspace{3mm}



\problem{\textbf{5}} Решите линейные диофантовы уравнения с двумя целыми неизвестными:

а.\textbf{(12)} $4x+7y=-11$

б.\textbf{(18)} $222156x-248292y=209088$



\problem{\textbf{6}}\textbf{(24)} Чему равны числа $A,B,C$, если верно следующее равенство:

$\dfrac{A}{x-1}-\dfrac{B}{x+2}+\dfrac{C}{3-x}=\dfrac{7x^2-10x+27}{(x-1)(x+2)(3-x)}$


\problem{\textbf{7}}\textbf{(20)} Пусть множество $E$ является пересечением множеств $A$ и $B$, а множество $M$ является
объединением множеств $C$ и $D$. Найдите пересечение и объединение множеств $E$ и $M$.

$A \in (-2026;-2002] \cup [-1971;-1953) \cup [-1945,0905;-1941] \cup \{-1917\} \cup (-1162;988]$ 

$B \in (-2012;-2002) \cup (-2000;-1941) \cup (-1924;-1914] \cup [-1161;2022,2402)$ 

$C \in  \{-2026,0702\} \cup \{-2025\} \cup [-2013;-476] \cup (67;5242] $

$D \in (-2012,0208;-456] \cup \{42\} \cup \{52\} \cup \{67\} \cup (1993;7101952]$ 


\problem{\textbf{8}}\textbf{(24)} АА купил своим ученикам 20 кг сникерсов, 30 кг чокопаек и 40 кг бэбифоксов, заплатив за всё 505 тысяч рублей. Если сникерсы подорожают на 30\%, чокопайки подешевеют на 20\%, бэбифоксы подорожают на 10\% , то за такую же массу соответствующих сладостей (20, 30, 40 кг) АА заплатит 526 тысяч рублей. Если сникерсы подешевеют на 40\%, чокопайки подорожают на 40\%, бэбифоксы подешевеют на 10\%, то за такую же массу соответствующих сладостей (20, 30, 40 кг) АА заплатит 507 тысяч рублей. Сколько рублей будет стоить набор, состоящий из 560 г сникерсов, 600 г чокопаек и 650 г бэбифоксов?

\problem{\textbf{9}}\textbf{([0;+$\infty))$} Если Вы ничего больше не можете решить, то в этом задании Вы можете написать все теоремы, признаки, свойства, интересные факты и т.д., которые связаны с математикой. 

\centerline{\textbf{Всего баллов 200$+$!}}
\centerline{\textbf{пэпэ}}